\chapter{Sfery Milnora: ten sam typ topologiczny,
 różne struktury gładkie}
\label{ch:sfery-Milnora}
\index{sfera Milnora@Sfera Milnora}

Twierdzenie~\ref{thm:poincare-high-24}
dało homeomorfizm sfer homotopijnych
z $S^n$. Nie twierdziło, że homeomorfizm
można wybrać gładki. Wymiar siedem
dostarcza konkretnego kontrprzykładu.
Zbudujemy go jako wiązkę sfer
$S^3\to M^7\to S^4$.
Występująca tutaj sfera $S^4$
jest \emph{bazą wiązki}; zbudowana
rozmaitość ma wymiar siedem.

\section{Sklejanie dwóch połówek wiązki}
\label{sec:milnor-clutching-25}

Utożsammy $\R^4$ z algebrą kwaternionów
$\mathbb H$ i sferę jednostkową
z grupą mnożenia $S^3\subset\mathbb H$.
Rozpiszmy $S^4$ na dwie półsfery
$D^4_+$ i $D^4_-$, sklejone wzdłuż
równika $S^3$. Nad każdą półsferą
weźmy trywialną wiązkę
$D^4_\pm\times\mathbb H$.
Na wspólnym brzegu identyfikujemy
włókna przez
\begin{equation}\label{eq:milnor-clutching-25}
 (u,v)_+\sim
 \bigl(u,u^hvu^j\bigr)_-,
 \qquad u\in S^3,\quad v\in\mathbb H,
 \quad h,j\in\Z.
\end{equation}
Mnożenie przez jednostkowy kwaternion
z lewej i prawej jest izometrią
orientacji dodatniej $\R^4$.
Przepis~\eqref{eq:milnor-clutching-25}
tworzy więc orientowaną wiązkę
rzeczywistych $4$-płaszczyzn
$E_{h,j}\to S^4$. Oznaczmy przez
$B_{h,j}=D(E_{h,j})$ jej wiązkę
dysków, a przez
$M_{h,j}=S(E_{h,j})=\partial B_{h,j}$
wiązkę sfer jednostkowych.
Przykładem jest $(h,j)=(1,0)$:
odpowiednia wiązka sfer jest
klasyczną fibracją Hopfa
$S^3\to S^7\to S^4$.

\begin{figure}[htbp]
\centering
\begin{tikzpicture}[>=Latex,font=\small,line cap=round]
 \shade[inner color=white,outer color=manifoldblue!43]
      (-2.65,0) circle (1.27);
 \shade[inner color=white,outer color=manifoldblue!43]
      (2.65,0) circle (1.27);
 \draw[manifoldblue!75!black,very thick]
      (-2.65,0) circle (1.27)
      (2.65,0) circle (1.27);
 \draw[manifoldblue!65!black,thick]
      (-3.92,0)--(-1.38,0)
      (1.38,0)--(3.92,0);
 \node at (-2.65,.55) {$D^4_+$};
 \node at (2.65,.55) {$D^4_-$};
 \node[align=center,font=\scriptsize]
      at (-2.65,-.55) {włókno $S^3$};
 \node[align=center,font=\scriptsize]
      at (2.65,-.55) {włókno $S^3$};
 \draw[->,very thick,manifoldgold!85!black]
      (-1.15,0)--(1.15,0);
 \node[above,align=center] at (0,.18)
      {$v\mapsto u^hvu^j$};
 \node[below] at (0,-.20)
      {$u\in S^3$};
\end{tikzpicture}
\caption{Schemat sklejania wiązki nad dwiema
półsferami $S^4$. Kolorowe okręgi oznaczają
półsfery w umownym przekroju; zmienna $u$
parametryzuje ich wspólny równik. Mapa na strzałce
skręca włókno i określa gładką wiązkę.}
\label{fig:milnor-clutching-25}
\end{figure}

\begin{stwierdzenie}[Dwie liczby charakterystyczne]
\label{prop:milnor-classes-25}
Niech $\alpha$ będzie dodatnim generatorem
$H^4(S^4;\Z)\cong\Z$. Przy spójnie
ustalonych orientacjach
\[
 e(E_{h,j})=(h+j)\alpha,\qquad
 p_1(E_{h,j})=\pm2(h-j)\alpha.
\]
Znak w drugim wzorze zależy od przyjętej
identyfikacji dwóch czynników $\operatorname{Spin}(4)$
z parametrami $(h,j)$. Przy ustalonych konwencjach
jest określony; w późniejszym rachunku
występuje jego kwadrat.
\end{stwierdzenie}
\begin{proof}
Klasa Eulera wiązki orientowanej
nad $S^4$ mierzy stopień
odwzorowania $S^3\to S^3$
otrzymanego z funkcji sklejającej
przez zastosowanie jej do jednego
ustalonego wektora jednostkowego.
Dla $v=1$ w~\eqref{eq:milnor-clutching-25}
jest to $u\mapsto u^{h+j}$.
Mnożenie punktowe odwzorowań
$S^3\to S^3$ dodaje ich klasy
w $\pi_3(S^3)\cong\Z$, więc stopień
tego odwzorowania wynosi $h+j$.
To dowodzi wzoru na $e$.

Wyprowadzimy również współczynnik przy $p_1$
z dwóch prostych wiązek. Dla klasy
charakterystycznej $c$ czwartego stopnia
(tutaj $c=e$ lub $c=p_1$)
przepis
\[
 [g]\in\pi_3(SO(4))\longmapsto
 \langle c(E_g),[S^4]\rangle\in\Z
\]
jest homomorfizmem. Rzeczywiście dodawanie
klas sklejania można przedstawić przez
odwzorowanie ściskające równik na dwie sfery;
fundamentalna klasa $S^4$ przechodzi wtedy
na sumę ich klas fundamentalnych, a
naturalność $c$ dodaje oba parowania.
Ponieważ lewa i prawa multiplikacja
komutują, klasa~\eqref{eq:milnor-clutching-25}
jest sumą $h$ kopii klasy $(1,0)$
i $j$ kopii klasy $(0,1)$.

Wiązka $E_{1,0}$ jest kwaternioniczną
wiązką Hopfa. Mnożenie z lewej strony
jest liniowe względem ustalonej struktury
zespolonej danej mnożeniem z prawej przez $i$.
Otrzymujemy zespoloną wiązkę rangi $2$
o grupie strukturalnej $SU(2)$, więc $c_1=0$.
Jej $c_2$ jest, z dokładnością do znaku
orientacji, klasą Eulera wiązki rzeczywistej,
czyli generatorem $H^4(S^4;\Z)$:
obie klasy liczą zera poprzecznego przekroju,
przy czym $c_2$ używa orientacji zespolonej.
Tożsamość
$p_1(V_{\R})=c_1(V)^2-2c_2(V)$
daje $p_1(E_{1,0})=\pm2\alpha$.
Można ją sprawdzić po rozszczepieniu
$V$ na linie z formalnymi klasami
$x,y$: kompleksyfikacja $V_{\R}$
ma pierwiastki $x,y,-x,-y$, zatem
$-c_2(V_{\R}\otimes\C)=x^2+y^2
=(x+y)^2-2xy$. To jest właśnie
definicja $p_1$ z
Rozdziału~\ref{ch:przeciecia-dualnosc-klasy}.

Sprzężenie kwaternionowe $v\mapsto\bar v$
przenosi przejście $v\mapsto vu$ dla
$E_{0,1}$ na $\bar v\mapsto u^{-1}\bar v$.
Jako wiązki bez ustalonej orientacji
$E_{0,1}$ i $E_{-1,0}$ są więc izomorficzne.
Klasa $p_1$ nie zależy od orientacji włókna,
stąd $p_1(E_{0,1})=-p_1(E_{1,0})$.
Addytywność z poprzedniego akapitu kończy
rachunek: $p_1(E_{h,j})=\pm2(h-j)\alpha$,
z jednym spójnym wyborem znaku.
Zasada rozszczepiania użyta w rachunku
klas jest przywołana także w
Sekcji~\ref{sec:chern-weil-26};
alternatywny rachunek przez
$\operatorname{Spin}(4)$ podaje lemat 3
\href{https://sites.math.rutgers.edu/~feehan/teaching/math866/milnor7sphere.pdf}
{Milnora, \emph{On Manifolds Homeomorphic to
the 7-Sphere} (1956)}.
\end{proof}

\section{Dlaczego całkowita przestrzeń jest
sferą homotopijną?}
\label{sec:milnor-homotopy-25}

Niech $h+j=1$. Zastosujmy
Mayera--Vietorisa do dwóch części
$M_{h,j}$ leżących nad półsferami.
Każda ma typ homotopijny $S^3$,
a ich część wspólna ma typ
$S^3_{\text{baza}}\times
S^3_{\text{włókno}}$.
Na grupach $H_3$ mapa z części
wspólnej do sumy homologii
obu połówek ma, po wyborze baz,
macierz
\begin{equation}\label{eq:milnor-MV-25}
 \begin{pmatrix}
 0&1\\
 h+j&1
 \end{pmatrix}.
\end{equation}
Pierwsza kolumna odpowiada
równikowi bazy. W pierwszej
trywializacji jest on brzegiem
dysku; w drugiej jego obraz
we włóknie ma stopień $h+j$.
Druga kolumna to sfera włókna,
generator $H_3$ w obu częściach.
Znaki jednego wiersza można
odwrócić zgodnie z konwencją
ciągu Mayera--Vietorisa;
wyznacznik ma wartość
bezwzględną $|h+j|=1$.
Z dokładności wynika zatem
$H_4(M_{h,j})=0$ i
$H_3(M_{h,j})=0$.
Pozostałe wyrazy tego samego
ciągu dają
\[
 H_i(M_{h,j};\Z)=
 \begin{cases}
 \Z,&i=0,7,\\
 0,&0<i<7.
 \end{cases}
\]
Rzeczywiście odwzorowanie łączące
$H_7(M_{h,j})\to H_6(S^3\times S^3)$
jest izomorfizmem, a w innych
stopniach dodatnich poza trzecim
grupy części są zerowe.

Długi ciąg homotopii fibracji
$S^3\to M_{h,j}\to S^4$ daje
$\pi_1(M_{h,j})=0$.
Wybierzmy osadzony zamknięty dysk
$D^7\subset M_{h,j}$. Zwińmy do punktu
zamknięty zbiór
$M_{h,j}\setminus\operatorname{int}D^7$.
Iloraz jest homeomorficzny z $S^7$,
otrzymujemy więc odwzorowanie
stopnia jeden $M_{h,j}\to S^7$.
Obliczone grupy homologii pokazują,
że jest ono izomorfizmem homologii.
Homologiczna wersja twierdzenia
Whiteheada
(Wniosek~\ref{cor:whitehead-homologiczny})
daje równoważność homotopijną.
Zatem $M_{h,j}$ jest gładką
sferą homotopijną, a przez
Twierdzenie~\ref{thm:poincare-high-24}
jest homeomorficzna z $S^7$.

\section{Niezmiennik wykrywający egzotyczność}
\label{sec:milnor-invariant-25}

Sama homologia i równoważność
homotopijna nie rozstrzygają
o dyfeomorfizmie. Potrzebujemy
niezmiennika gładkiego.
Załóżmy, że $M^7$ jest
zorientowaną rozmaitością
z $H^3(M;\Z)=H^4(M;\Z)=0$
i że $M=\partial B$ dla
zorientowanej gładkiej
$8$-rozmaitości $B$.
W długim ciągu pary
odwzorowanie
$j:H^4(B,M;\Z)\to H^4(B;\Z)$
jest izomorfizmem.
Dlatego $p_1(TB)$ ma
jednoznaczne względne podniesienie
$\bar p_1=j^{-1}p_1(TB)$.
Połóżmy
\[
 q(B)=\langle
    \bar p_1\smile p_1(TB),[B,M]\rangle,
 \qquad
 \lambda(M)=[\,2q(B)-\operatorname{sign}(B)\,]
                    \in\Z/7.
\]
Sygnatura $\operatorname{sign}(B)$
jest sygnaturą formy przecięcia
na obrazie
$H_4(B)\to H_4(B,M)$.
Na przykład dla $M=S^7$
można wziąć $B=D^8$,
co daje $q(B)=0$ i
$\lambda(S^7)=0$.

\begin{twierdzenie}[Milnor: niezależność
od wypełnienia]
\label{thm:milnor-invariant-25}
Powyższa klasa $\lambda(M)\in\Z/7$
nie zależy od wybranego gładkiego
zorientowanego wypełnienia $B$.
Jest więc niezmiennikiem
zorientowanego dyfeomorfizmu $M$.
\end{twierdzenie}
\begin{proof}
Niech $B_0,B_1$ będą dwoma
wypełnieniami. Sklejmy
$C=B_0\cup_M(-B_1)$;
to zamknięta zorientowana
$8$-rozmaitość. Z przyjętych założeń
mamy $H^3(M)=H^4(M)=0$.
Długie ciągi par $(B_i,M)$ dają
izomorfizmy
$H^4(B_i,M)\xrightarrow{\sim}H^4(B_i)$.
W ciągu Mayera--Vietorisa dla $C$
znikają $H^3(M)$ i $H^4(M)$,
więc klasy czwarte obu części
sklejają się jednoznacznie. Względne
klasy $\bar p_1$ ograniczają się
do obu kawałków, a parowanie z $[C]$
jest sumą parowań z $[B_0,M]$
i $-[B_1,M]$. Tak samo forma przecięcia
na $C$ jest sumą ortogonalną form
na kawałkach, z odwróconym znakiem
na drugim. Dlatego
\[
 q(C)=q(B_0)-q(B_1),\qquad
 \operatorname{sign}(C)
  =\operatorname{sign}(B_0)
   -\operatorname{sign}(B_1).
\]
Na zamkniętej $C$ mamy
$q(C)=\langle p_1(C)^2,[C]\rangle$.
Zastosujmy wzór Hirzebrucha
na sygnaturę w wymiarze osiem:
\[
 45\operatorname{sign}(C)=
 \langle 7p_2(C)-p_1(C)^2,[C]\rangle.
\]
Współczynniki $7,-1,45$ można sprawdzić
formalnie bez dowodzenia samego twierdzenia
Hirzebrucha. Dla formalnych pierwiastków
Pontriagina $x_a$ rozwijamy
$x/\tanh x=1+x^2/3-x^4/45+\cdots$.
Składnik stopnia osiem iloczynu wynosi
\[
 \frac1{9}\sum_{a<b}x_a^2x_b^2
 -\frac1{45}\sum_a x_a^4
 =\frac{7p_2-p_1^2}{45},
 \quad p_1=\sum_a x_a^2,\quad
 p_2=\sum_{a<b}x_a^2x_b^2.
\]
Głębokim wejściem jest identyfikacja
parowania tego wielomianu z sygnaturą
każdej \emph{zamkniętej} zorientowanej
$8$-rozmaitości; nie stosujemy jej
bezpośrednio do $B$ z brzegiem. W tej
postaci używa jej
\href{https://sites.math.rutgers.edu/~feehan/teaching/math866/milnor7sphere.pdf}
{Milnor, \emph{On Manifolds Homeomorphic
to the 7-Sphere}, §1}.
Stąd
\[
 45\bigl(2q(C)-\operatorname{sign}(C)\bigr)
 =
 \langle 91p_1(C)^2-7p_2(C),[C]\rangle.
\]
Prawa strona dzieli się przez $7$.
Ponieważ $45$ jest odwracalne
modulo $7$, również
$2q(C)-\operatorname{sign}(C)$
jest podzielne przez $7$.
Po podstawieniu wzorów
addytywności otrzymujemy
$\lambda_{B_0}(M)=\lambda_{B_1}(M)$.
Użyty wzór Hirzebrucha jest
głębokim wynikiem zewnętrznym;
to właśnie on jest istotną
częścią argumentu Milnora.
\end{proof}

\begin{stwierdzenie}[Rachunek dla wiązek Milnora]
\label{prop:milnor-value-25}
Dla nieparzystego $k$ wybierzmy
\[
 h=\frac{k+1}{2},\qquad
 j=\frac{1-k}{2},
 \qquad M_k=M_{h,j}.
\]
Wtedy
\[
 \lambda(M_k)
 =8k^2-1\equiv k^2-1\pmod7.
\]
\end{stwierdzenie}
\begin{proof}
Mamy $h+j=1$, więc
poprzednia sekcja pokazuje,
że $M_k$ jest sferą homotopijną.
Wiązka dysków $B_k=D(E_{h,j})$
deformuje się na zerowy przekrój
$S^4$, dlatego
$H^4(B_k;\Z)=\Z\alpha$.
Przez Stwierdzenie~\ref{prop:milnor-classes-25}
oraz stabilne rozszczepienie
$TB_k\cong\pi^*TS^4\oplus\pi^*E_{h,j}$
mamy
$p_1(TB_k)=\pm2k\alpha$:
wiązka $TS^4$ jest stabilnie
trywialna, zatem jej $p_1=0$.
Forma przecięcia w środku
$B_k$ ma macierz $[1]$,
ponieważ samoprzecięcie
zerowego przekroju jest liczbą
Eulera $e(E_{h,j})=1$.
Stąd $\operatorname{sign}(B_k)=1$.
Twierdzenie Thoma
(Twierdzenie~\ref{thm:euler-thom-16})
daje generator $U\in H^4(B_k,M_k;\Z)$
oraz tożsamość $j(U)=\pi^*e(E_{h,j})$.
Stąd $U\smile j(U)=
\pi^*e(E_{h,j})\smile U$, a parowanie
tej klasy z $[B_k,M_k]$ jest równe
$\langle e(E_{h,j}),[S^4]\rangle=1$.
Otrzymujemy
$q(B_k)=(2k)^2=4k^2$,
a po podstawieniu do definicji
$\lambda$ dostajemy
$2\cdot4k^2-1$.
\end{proof}

\begin{twierdzenie}[Egzotyczna sfera Milnora]
\label{thm:exotic-milnor-25}
Rozmaitość $M_3^7$ otrzymana
dla $(h,j)=(2,-1)$
jest homeomorficzna z $S^7$,
ale nie jest z nią dyfeomorficzna.
\end{twierdzenie}
\begin{proof}
Warunek $h+j=1$ i obliczenie
homologii pokazują, że
$M_3\simeq S^7$; homeomorfizm
wynika z
Twierdzenia~\ref{thm:poincare-high-24}.
Natomiast
$\lambda(M_3)=8\cdot9-1=71
\equiv1\pmod7$, podczas gdy
$\lambda(S^7)=0$.
Gdyby istniał dyfeomorfizm
zachowujący orientację, te
niezmienniki musiałyby być równe.
Odwrócenie orientacji zmienia
znak $\lambda$, lecz $\pm1$
nadal nie jest zerem modulo $7$.
Zatem nie istnieje także
dyfeomorfizm odwracający
orientację.
\end{proof}

W języku
Definicji~\ref{def:structure-set-23}
przykład daje strukturę na
typie homotopijnym $S^7$
różną od standardowej.
Nie oznacza to, że każdą
różnicę struktur wykrywa
$\lambda$; obliczenie pełnej
grupy sfer homotopijnych
$\Theta_7\cong\Z/28$ jest
osobnym wynikiem
\href{https://www.sas.rochester.edu/mth/sites/doug-ravenel/otherpapers/kervaire-milnor.pdf}
{M.\ Kervaire'a i J.\ Milnora,
\emph{Groups of Homotopy Spheres I}
(1963)}. Tutaj wystarcza
jedna jawnie skonstruowana
egzotyczna sfera.

\par\smallskip
{\footnotesize\noindent\emph{Dalsza lektura:}
\href{https://sites.math.rutgers.edu/~feehan/teaching/math866/milnor7sphere.pdf}
{J.\ Milnor, \emph{On Manifolds Homeomorphic
to the 7-Sphere}, Ann. of Math. 64 (1956),
§1--3}. W oryginale homeomorfizm
uzyskano przez specjalną funkcję
Morse'a; tutaj użyliśmy
Twierdzenia~\ref{thm:poincare-high-24}.\par}

\clearpage
